\documentclass[11pt]{article}

\usepackage[T1]{fontenc}
\usepackage[utf8]{inputenc}
\usepackage[french]{babel}
\usepackage[margin=1in]{geometry}
\usepackage{amsmath,amssymb,amsthm}
\usepackage[hidelinks]{hyperref}
\usepackage{listings}
\usepackage{booktabs}
\usepackage{xcolor}

\lstset{basicstyle=\ttfamily\small,breaklines=true,columns=fullflexible}
\setlength{\emergencystretch}{3em}

\newtheorem{theorem}{Théorème}
\newtheorem{lemma}[theorem]{Lemme}
\newtheorem{corollary}[theorem]{Corollaire}
\newcommand{\NN}{\mathbb{N}}
\newcommand{\QQ}{\mathbb{Q}}
\newcommand{\bind}[2]{\binom{#1}{#2}}
\newcommand{\wedgep}[1]{\textstyle\bigwedge^{#1}}

\title{Deux techniques de preuve pour le théorème des deux familles :\\
permutations, puissances extérieures et position générale dans Lean 4}
\author{Louis Gestin\\\small Élève indépendant, France\\\small \href{mailto:louis@gestin.net}{louis@gestin.net}}
\date{Rapport technique indépendant, juillet 2026}

\begin{document}
\maketitle

\begin{abstract}
Nous présentons une formalisation Lean~4 / Mathlib de la famille de théorèmes de
Bollobás--Lovász--Frankl sur les deux familles, développée selon deux techniques de
preuve mathématiquement distinctes. La première est une preuve par dénombrement de
permutations de l'inégalité pondérée de Bollobás pour les paires d'ensembles
~\cite{bollobas1965}, ainsi que de son corollaire uniforme. La seconde est une preuve par
algèbre extérieure du théorème asymétrique de Lovász--Frankl pour les sous-espaces
~\cite{lovasz1977}, reliée au théorème asymétrique de Frankl--Kalai pour les ensembles
finis~\cite{frankl1982,kalai1984} par une construction en position générale fondée sur
la courbe des moments sur~$\QQ$. Les cinq théorèmes principaux sont entièrement vérifiés
par machine et ne dépendent que des axiomes standards \texttt{propext},
\texttt{Classical.choice} et \texttt{Quot.sound}; le développement ne contient aucun
\texttt{sorry}, \texttt{admit} ni axiome personnalisé. Le code comprend environ $1440$
lignes réparties dans $13$ fichiers ($51$ déclarations, dont $37$ théorèmes) et se
compile en environ $133$~s lorsque le cache Mathlib épinglé est déjà disponible. Notre
thèse est que, bien que les deux preuves utilisent des bibliothèques différentes, elles
produisent le même objet abstrait : un certificat triangulaire d'indépendance. Nous
isolons plusieurs composants réutilisables pour Mathlib : un principe d'indépendance
linéaire triangulaire, des lemmes d'annulation, de non-annulation et de concaténation
pour les vecteurs extérieurs décomposables, un dénombrement exact sans division des
événements de séparation, ainsi qu'un théorème de position générale pour la courbe des
moments et ses conséquences sur les intersections d'espaces engendrés. Ce travail est
une formalisation et n'introduit pas de nouveau résultat mathématique.
\end{abstract}

\section{Introduction}\label{sec:intro}

Les théorèmes des deux familles sont des résultats fondamentaux de combinatoire
extrémale. Dans leur forme uniforme la plus simple, ils affirment que, si
$A_1,\dots,A_m$ sont des ensembles de taille $a$ et $B_1,\dots,B_m$ des ensembles de
taille $b$, avec $A_i\cap B_i=\varnothing$ et $A_i\cap B_j\neq\varnothing$ pour
$i\neq j$, alors $m\le\bind{a+b}{a}$. La version pondérée de
Bollobás~\cite{bollobas1965} supprime l'hypothèse d'uniformité et donne la borne
rationnelle exacte
$\sum_i 1/\bind{|A_i|+|B_i|}{|A_i|}\le 1$. Lovász~\cite{lovasz1977} a remplacé les
ensembles par des sous-espaces et démontré une version \emph{asymétrique}, où
l'intersection croisée n'est exigée que pour $i<j$; sa preuve naturelle utilise
l'algèbre extérieure. Frankl~\cite{frankl1982} et Kalai~\cite{kalai1984} ont ensuite
transporté cette borne asymétrique vers les ensembles finis grâce à une représentation
en position générale. Une présentation pédagogique apparaît dans les notes de
Babai--Frankl~\cite{babaifrankl1992}.

Ces preuves mettent en jeu deux techniques réellement différentes. L'inégalité pondérée
se démontre en comptant, parmi tous les ordres linéaires de l'ensemble de base, les
« événements de séparation » dans lesquels chaque $A_i$ apparaît entièrement avant
$B_i$; la disjonction de ces événements pour des indices différents donne la borne. Le
théorème sur les sous-espaces associe à chaque paire un vecteur extérieur décomposable et
une application de multiplication, puis applique un argument triangulaire
d'indépendance linéaire dans $\wedgep{a}V$, de dimension exactement
$\bind{a+b}{a}$. Nous formalisons les deux preuves dans Lean~4~\cite{lean4}, au-dessus de
Mathlib~\cite{mathlib2020}.

\paragraph{Thèse.} Notre observation centrale est que les deux preuves, malgré
l'utilisation de parties disjointes de la bibliothèque, convergent vers le même type de
certificat abstrait : une famille d'objets « non nuls sur la diagonale et nuls strictement
sous celle-ci ». Dans la preuve combinatoire, il s'agit d'événements de séparation deux à
deux disjoints et de taille exactement calculée. Dans la preuve algébrique, il s'agit
d'applications tests de multiplication extérieure qui alimentent un unique lemme
générique, \lstinline{LinearIndependent.of_upperTriangular_maps}. La formalisation rend
cette correspondance précise plutôt que simplement métaphorique.

\paragraph{Contributions.}
\begin{enumerate}
\item Un développement Lean complet, vérifié par le noyau, des théorèmes pondéré et
  uniforme de Bollobás, du théorème de Lovász--Frankl pour les sous-espaces, du théorème
  de position générale de la courbe des moments et du théorème asymétrique de
  Frankl--Kalai pour les ensembles, selon deux voies de preuve indépendantes
  (sections~\ref{sec:count}--\ref{sec:bridge}).
\item Une infrastructure réutilisable pour laquelle nous n'avons pas trouvé
d'équivalent direct dans Mathlib v4.28.0 : indépendance triangulaire obtenue par des
applications tests, non-annulation et concaténation des décomposables extérieurs,
dénombrement exact et sans division des énumérations séparatrices, et position générale
de la courbe des moments avec des lemmes d'intersection des espaces engendrés
(section~\ref{sec:reuse}).
\item Un audit sémantique indépendant, des tests adversariaux des hypothèses et une
évaluation mesurée (section~\ref{sec:eval}), permettant des affirmations prudentes sur
l'originalité de la formalisation (section~\ref{sec:related}).
\end{enumerate}

\section{Vue d'ensemble mathématique}\label{sec:overview}

Dans toute la suite, $\alpha$ est un type de base fini et $\iota$ un type d'indices fini;
$n=|\alpha|$.

\begin{theorem}[Bollobás pondéré]\label{thm:weighted}
Pour des ensembles $A_i,B_i\subseteq\alpha$ tels que $A_i\cap B_i=\varnothing$ pour
tout $i$ et $A_i\cap B_j\neq\varnothing$ pour tous $i\neq j$,
\[ \sum_i \frac{1}{\bind{|A_i|+|B_i|}{|A_i|}}\le 1. \]
\end{theorem}

\begin{corollary}[Bollobás uniforme]\label{thm:uniform}
Si, de plus, $|A_i|=a$ et $|B_i|=b$ pour tout $i$, alors
$|\iota|\le\bind{a+b}{a}$. La borne est atteinte par la famille des complémentaires sur
$\{1,\dots,a+b\}$.
\end{corollary}

\begin{theorem}[Lovász--Frankl, sous-espaces]\label{thm:subspace}
Soit $\dim V=a+b$, et soient $U_i,W_i\le V$ tels que $\dim U_i=a$, $\dim W_i=b$,
$U_i\cap W_i=0$ et $U_i\cap W_j\neq0$ pour $i<j$. Alors
$m\le\bind{a+b}{a}$.
\end{theorem}

\begin{theorem}[Frankl--Kalai, ensembles]\label{thm:fk}
Soient $A_i,B_i\subseteq\alpha$ avec $|A_i|=a$, $|B_i|=b$,
$A_i\cap B_i=\varnothing$ et $A_i\cap B_j\neq\varnothing$ pour $i<j$. Alors
$m\le\bind{a+b}{a}$.
\end{theorem}

\noindent\emph{Esquisses de preuve.} Pour le théorème~\ref{thm:weighted}, un ordre
linéaire uniforme de $\alpha$ sépare $(A_i,B_i)$ avec une « probabilité »
$1/\bind{|A_i|+|B_i|}{|A_i|}$. Les événements de séparation associés à deux indices
distincts sont disjoints, car une paire croisée imposerait deux inégalités d'ordre
contradictoires; la somme des inverses est donc au plus $1$. Nous formalisons cet
argument par un dénombrement entier exact plutôt que par une probabilité. Pour le
théorème~\ref{thm:subspace}, les vecteurs décomposables
$u_i=\bigwedge U_i\in\wedgep{a}V$ et la multiplication par
$w_j=\bigwedge W_j$ donnent des applications telles que $u_i\wedge w_i\neq0$ sur la
diagonale et $u_i\wedge w_j=0$ pour $i<j$. La triangularité implique l'indépendance des
$u_i$, puis $m\le\dim\wedgep{a}V=\bind{a+b}{a}$. Le théorème~\ref{thm:fk} suit en
réalisant les ensembles comme des espaces engendrés par des vecteurs de la courbe des
moments, de sorte que la disjonction devient une somme directe et qu'une intersection
non vide fournit un vecteur commun non nul.

\section{Architecture de la formalisation}\label{sec:arch}

Le développement est modulaire et se divise en deux voies qui ne partagent que l'API de
prédicats \lstinline{SetPairs.Basic}. La figure~\ref{fig:dep} montre le graphe acyclique
des imports locaux, construit à partir des lignes
\lstinline{import RequestProject.*} et vérifié dans
\texttt{DEPENDENCY\_AUDIT.md}.

\begin{figure}[t]
\centering
\begin{verbatim}
                         SetPairs/Basic
                         /            \
          PermutationOrders          SkewSetPairs
                |                      /        \
          SeparationEvents           /          \
                |             MomentCurve   TwoFamiliesSubspaces
          WeightedBollobas                 /     |         \
                |                Triangular  ExteriorDecomp  ExteriorMult
             Uniform                        Independence
                \                              /
                 \____________ Main _________/    (+ Examples)
\end{verbatim}
\caption{Graphe des dépendances entre modules. Branche gauche : dénombrement de
permutations. Branche droite : puissances extérieures et position générale. Les deux
branches ne se rejoignent que dans \texttt{SetPairs/Basic} et dans \texttt{Main}.}
\label{fig:dep}
\end{figure}

\begin{itemize}
\item \texttt{SetPairs/} : \texttt{Basic} (prédicats), \texttt{PermutationOrders},
  \texttt{SeparationEvents}, \texttt{WeightedBollobas}, \texttt{Uniform}.
\item \texttt{LinearAlgebra/} : \texttt{TriangularIndependence},
  \texttt{ExteriorDecomposable}, \texttt{ExteriorMultiplication},
  \texttt{TwoFamiliesSubspaces}, \texttt{MomentCurve}.
\item \texttt{SkewSetPairs}, \texttt{Examples}, \texttt{Main}.
\end{itemize}

\section{Dénombrement exact des permutations}\label{sec:count}

Une \emph{énumération ordonnée} de $\alpha$ est une équivalence
$e:\alpha\simeq\mathrm{Fin}\,n$; il en existe $n!$
(\lstinline{card_orderEnumerations}, conséquence de
\lstinline{Fintype.card_equiv}). La paire $(A,B)$ est \emph{séparée} par $e$ lorsque
tout élément de $A$ précède tout élément de $B$ :
\begin{verbatim}
def Separates (e : OrderEnumeration a) (A B : Finset a) : Prop :=
  forall {x}, x in A -> forall {y}, y in B -> e x < e y
\end{verbatim}
et $\mathrm{separationEvent}(A,B)$ est le \texttt{Finset} des énumérations qui
satisfont cette propriété.

\paragraph{Dénombrement exact.} Le coeur de l'argument est l'identité
\[ |\mathrm{separationEvent}(A,B)|\cdot\bind{|A|+|B|}{|A|}=n!, \]
obtenue à partir de la formule fermée
(\lstinline{card_separationEvent_eq})
\[ |\mathrm{separationEvent}(A,B)|=\bind{n}{a+b}\,a!\,b!\,(n-(a{+}b))!, \]
où $a=|A|$, $b=|B|$ et $S=A\cup B$. La preuve repose sur une véritable équivalence de
factorisation des ordres : une énumération séparatrice est déterminée par trois données
indépendantes :
\begin{itemize}
\item l'ensemble de positions $P=e(S)\subseteq\mathrm{Fin}\,n$, de taille $a+b$, qui
  fournit le facteur $\bind{n}{a+b}$;
\item l'ordre relatif de $S$ dans $P$, avec tous les éléments de $A$ avant ceux de $B$,
  qui fournit $a!\,b!$;
\item l'ordre relatif de $S^{\mathsf c}$ dans $P^{\mathsf c}$, qui fournit
  $(n-a-b)!$.
\end{itemize}
Plus précisément :
\begin{itemize}
\item \lstinline{card_separates_full} : lorsque $A\cup B=\mathrm{univ}$, les
  énumérations séparatrices sont en bijection avec
  $(A\simeq\mathrm{Fin}\,a)\times(B\simeq\mathrm{Fin}\,b)$, d'où $a!\,b!$. La
  construction utilise deux injections explicites et les faits d'ordre
  \lstinline{separates_val_lt} et \lstinline{card_le_separates_val}.
\item \lstinline{card_separates_toP} transporte ce dénombrement le long de
  l'isomorphisme ordonné $P\simeq_o\mathrm{Fin}(a+b)$.
\item \lstinline{fiber_card_eq} identifie, pour un $P$ fixé, la fibre des énumérations
  vérifiant $e(S)=P$ avec le produit de l'ordre relatif de $S$ et d'une équivalence entre
  les complémentaires.
\item \lstinline{card_separationEvent_eq} somme les tailles de ces fibres sur les
  $\bind{n}{a+b}$ ensembles de positions, via
  \lstinline{Finset.card_eq_sum_card_fiberwise}.
\end{itemize}
L'identité factorielle finale est obtenue \emph{sans division dans les naturels}, à
partir de $\bind{a+b}{a}a!b!=(a+b)!$ et de
$\bind{n}{a+b}(a+b)!(n-(a+b))!=n!$, en appliquant deux fois
\lstinline{Nat.choose_mul_factorial_mul_factorial}.

\paragraph{Pourquoi ne pas utiliser les probabilités ?} Une formulation probabiliste
directe introduirait des mesures à valeurs dans $\mathbb{R}$ et des divisions, dont la
gestion formelle serait plus lourde que l'identité entière. Elle masquerait également
la structure des fibres, qui constitue le véritable coeur de la preuve. Le dénombrement
entier peut en outre être contrôlé exactement par les tests de régression
(section~\ref{sec:eval}).

\paragraph{Disjonction et assemblage.} Les événements de séparation sont deux à deux
disjoints. Si $x\in A_i\cap B_j$ et $y\in A_j\cap B_i$, séparer $(A_i,B_i)$ impose
$e(x)<e(y)$, tandis que séparer $(A_j,B_j)$ impose $e(y)<e(x)$
(\lstinline{separationEvents_disjoint}). C'est exactement à cet endroit que
l'intersection croisée dans \emph{les deux directions} est utilisée. La somme des
taille des événements disjoints satisfait $\sum_i|E_i|\le n!$; comme
$1/\bind{|A_i|+|B_i|}{|A_i|}=|E_i|/n!$, on obtient
\lstinline{weighted_bollobas}. Le corollaire uniforme
\lstinline{uniform_bollobas} met en facteur l'inverse constant, et
\lstinline{uniform_bollobas_sharp} formalise la famille des complémentaires qui atteint
la borne.

\section{Indépendance triangulaire et produits extérieurs}\label{sec:ext}

\paragraph{Indépendance triangulaire.} Si des applications linéaires $T_j$ satisfont
$T_j(u_j)\neq0$ et $T_j(u_i)=0$ pour $i<j$, alors les $u_i$ sont linéairement
indépendants (\lstinline{LinearIndependent.of_upperTriangular_maps}) :
\begin{verbatim}
theorem LinearIndependent.of_upperTriangular_maps
    (u : Fin m -> M) (T : Fin m -> M ->L[K] N)
    (hdiag : forall j, T j (u j) != 0)
    (hzero : forall {i j}, i < j -> T j (u i) = 0) :
    LinearIndependent K u
\end{verbatim}
La preuve considère une combinaison nulle $\sum_i g_i u_i=0$, choisit le plus grand
indice $j$ tel que $g_j\neq0$, puis applique $T_j$. Les termes d'indice $i<j$ sont
annulés par \texttt{hzero}, et ceux d'indice $i>j$ ont un coefficient nul par maximalité.
Il reste $g_jT_j(u_j)=0$, contradiction puisque $K$ est un corps. Le choix de l'indice
maximal correspond à l'orientation triangulaire supérieure; le dual
\lstinline{of_lowerTriangular_maps} utilise l'indice minimal.

\paragraph{Vecteurs décomposables.} Le produit alterné
$\iota\mathrm{Multi}\,K\,r\,v$ s'annule exactement lorsque $v$ est linéairement
dépendante. Le sens « dépendance implique annulation » est fourni par
\lstinline{AlternatingMap.map_linearDependent}; le sens de non-annulation
(\texttt{\(\iota\)Multi\_ne\_zero\_of\_linearIndependent}, dans la puissance extérieure
et dans l'algèbre extérieure) est déduit du lemme Mathlib
\texttt{\(\iota\)Multi\_family\_linearIndependent\_field}. L'identité de concaténation
$\iota\mathrm{Multi}\,a\,u\cdot\iota\mathrm{Multi}\,b\,w=
\iota\mathrm{Multi}\,(a+b)\,(u\mathbin{+\!\!+}w)$
(\texttt{\(\iota\)Multi\_mul\_\(\iota\)Multi}) transforme le produit de deux
décomposables en le décomposable de la famille concaténée. Sa preuve repose sur une
identité élément par élément entre \lstinline{List.ofFn} et \lstinline{Fin.append}.

\paragraph{Théorème sur les sous-espaces.} Pour chaque $i$, nous choisissons des bases de
$U_i$ et $W_i$ à l'aide de \lstinline{Module.finBasisOfFinrankEq}, en utilisant
$\dim U_i=a$ et $\dim W_i=b$. On forme $u_i\in\wedgep{a}V$ et une application test
$T_j=(\,\cdot\,*w_j)\circ\iota$, où $w_j$ est le décomposable associé à $W_j$.
Alors $T_j(u_i)$ est le décomposable des bases concaténées. Il est non nul sur la
diagonale, car $U_i\cap W_i=0$ rend la famille réunie indépendante
(\lstinline{append_wedge_ne_zero}), et nul pour $i<j$, car
$U_i\cap W_j\neq0$ la rend dépendante (\lstinline{append_wedge_eq_zero}). Le principe
triangulaire donne
$m\le\dim\wedgep{a}V=\bind{a+b}{a}$ par
\lstinline{exteriorPower.finrank_eq}. Les cas dégénérés $a=0$, $b=0$ et $m=0$ ne
nécessitent aucune séparation : les bases vides et
$\iota\mathrm{Multi}\,K\,0=1$ sont traitées uniformément par l'algèbre extérieure.

\section{Passage par la courbe des moments}\label{sec:bridge}

Le vecteur de la courbe des moments est
$v_x=(1,t_x,\dots,t_x^{d-1})\in\QQ^d$, où
$t:\alpha\hookrightarrow\QQ$ est injective. Toute famille d'au plus $d$ vecteurs
distincts de cette forme est linéairement indépendante
(\lstinline{momentCurve_linearIndependent}). La preuve utilise l'interpolation
polynomiale : pour isoler l'indice $i$, on prend
$\prod_{j\neq i}(X-t_j)$, de degré $|s|-1<d$. Ce choix évite le développement du
déterminant de Vandermonde et traite directement le cas \emph{rectangulaire} $|s|<d$,
sans réduction artificielle à une matrice carrée.

On en déduit
$\dim\operatorname{span}\{v_x:x\in S\}=|S|$ lorsque $|S|\le d$
(\lstinline{finrank_spanMomentCurve}), la trivialité de
$\operatorname{span}(S)\cap\operatorname{span}(T)$ lorsque $S$ et $T$ sont disjoints et
$|S|+|T|\le d$ (\lstinline{inf_spanMomentCurve_eq_bot}), et sa non-trivialité lorsque
$S\cap T\neq\varnothing$ (\lstinline{inf_spanMomentCurve_ne_bot}). Dans ce dernier cas,
un élément commun fournit un vecteur de la courbe des moments non nul, sa coordonnée
d'indice zéro valant $1$. Le théorème de Frankl--Kalai injecte ensuite le type de base
dans $\QQ$, via \lstinline{Fintype.equivFin} puis l'inclusion
$\NN\hookrightarrow\QQ$, représente chaque ensemble par son espace engendré dans
$\QQ^{a+b}$ et applique le théorème sur les sous-espaces. Le type de base peut être
strictement plus grand que $a+b$; seule la dimension ambiante $d=a+b$ est contrainte.

\section{Infrastructure réutilisable et lacunes de Mathlib}\label{sec:reuse}

Le tableau~\ref{tab:correspondence} associe les objets mathématiques centraux à leur
représentation Lean. Plusieurs composants sont génériques et n'ont pas été trouvés dans
Mathlib v4.28.0 lors de notre recherche documentée
(\texttt{MATHLIB\_PR\_PLAN.md}) : indépendance triangulaire par applications tests,
non-annulation et concaténation des décomposables extérieurs, dénombrement exact des
événements de séparation, et position générale de la courbe des moments avec ses lemmes
d'intersection. Mathlib fournit déjà l'API des puissances extérieures, leur dimension
$\dim\wedgep{a}V=\bind{\dim V}{a}$, le déterminant de Vandermonde et l'inégalité LYM,
qui correspond à une autre technique de dénombrement par chaînes.

\begin{table}[t]
\centering\scriptsize
\begin{tabular}{@{}llll@{}}
\toprule
Objet mathématique & Représentation Lean & Déclaration clé & Difficulté principale \\
\midrule
Ordre total & $\alpha\simeq\mathrm{Fin}\,n$ & \texttt{OrderEnumeration} & compter les bijections \\
Événement de séparation & \texttt{Finset} filtré & \texttt{separationEvent} & taille exacte des fibres \\
Décomposable extérieur & $\iota$\texttt{Multi}$\,K\,r\,v$ & \texttt{\(\iota\)Multi\_ne\_zero\_*} & non-annulation \\
Test par multiplication & \texttt{mulRight} $\circ$ inclusion & \texttt{\(\iota\)Multi\_mul\_\(\iota\)Multi} & concaténation \\
Vecteur moment & $\mathrm{Fin}\,d\to\QQ$ & \texttt{momentCurve} & indépendance rectangulaire \\
Espace engendré & \texttt{Submodule.span} & \texttt{spanMomentCurve} & intersection / somme directe \\
Intersection asymétrique & $i<j\to(A_i\cap B_j)\neq\varnothing$ & \texttt{SkewCrossIntersecting} & orientation \\
\bottomrule
\end{tabular}
\caption{Correspondance entre la preuve informelle et Lean.}
\label{tab:correspondence}
\end{table}

\begin{table}[t]
\centering\scriptsize
\begin{tabular}{@{}lllccl@{}}
\toprule
Résultat & Méthode & Déclaration Lean & Vérifié & Optimal & Dépendances \\
\midrule
Pondéré & dénombrement & \texttt{weighted\_bollobas} & oui & --- & compte séparateur \\
Uniforme & dénombrement & \texttt{uniform\_bollobas} & oui & oui & pondéré \\
Triangulaire & alg. linéaire & \texttt{of\_upperTriangular\_maps} & oui & --- & générique \\
Sous-espaces & extérieur & \texttt{lovasz\_frankl\_subspaces} & oui & --- & triangulaire, extérieur \\
Courbe des moments & interpolation & \texttt{momentCurve\_linearIndependent} & oui & --- & générique \\
Frankl--Kalai & passage & \texttt{frankl\_kalai\_skew} & oui & oui & sous-espaces, moments \\
\bottomrule
\end{tabular}
\caption{Couverture des théorèmes. « Optimal » signifie qu'une famille atteignant la
borne est formalisée; la construction par complémentaires convient aux bornes uniforme
et asymétrique.}
\label{tab:coverage}
\end{table}

\section{Évaluation}\label{sec:eval}

\paragraph{Versions.} Lean \texttt{v4.28.0}; Mathlib épinglé à
\texttt{v4.28.0}, commit \texttt{8f9d9cff}. Une compilation propre du projet avec les
artefacts Mathlib déjà présents réussit en $8039$ tâches et environ $133$~s dans notre
environnement. Une reconstruction froide doit d'abord télécharger les dépendances
épinglées et récupérer le cache Mathlib; elle est donc nettement plus longue.

\paragraph{Taille.} Le développement contient $13$ fichiers Lean, $1438$ lignes au
total ($1175$ non vides) et $51$ déclarations : $37$ théorèmes ou lemmes, $10$
définitions, $2$ abréviations et $2$ instances. Le module le plus long est
\texttt{SeparationEvents.lean} avec $481$ lignes, devant
\texttt{MomentCurve.lean} ($179$) et \texttt{TwoFamiliesSubspaces.lean} ($144$).

\paragraph{Axiomes.} Pour chaque théorème principal, \lstinline{#print axioms} affiche
exactement \lstinline{[propext, Classical.choice, Quot.sound]}, sans
\lstinline{sorryAx}. Le dépôt ne contient aucun \lstinline{sorry}, \lstinline{admit},
\lstinline{axiom} personnalisé, \lstinline{unsafe} ou
\lstinline{@[implemented_by]}, comme le vérifie
\texttt{scripts/check\_forbidden.sh}.

\paragraph{Tests de régression.} Le script
\texttt{scripts/verify\_small\_cases.py} vérifie en arithmétique rationnelle exacte
l'inégalité pondérée pour des ensembles de base de taille $2$ à $5$, la borne uniforme
pour les paramètres admissibles jusqu'à une taille de base $6$, la borne asymétrique
jusqu'à une taille $4$, ainsi que la construction optimale. Le script
\texttt{scripts/adversarial\_checks.py} fournit des contre-exemples exacts lorsque les
hypothèses principales sont retirées. Des tests Lean avec \lstinline{decide} et
\lstinline{norm_num} contrôlent également des valeurs binomiales, cinq cas limites du
dénombrement fermé et deux sommes adversariales strictement supérieures à $1$. Ces
calculs servent de tests de régression et ne remplacent pas les preuves vérifiées par le
noyau.

\paragraph{Composants réutilisables.} L'infrastructure générique des
sections~\ref{sec:ext}--\ref{sec:reuse} -- indépendance triangulaire, résultats sur les
décomposables extérieurs, position générale de la courbe des moments et dénombrement des
événements de séparation -- constitue l'essentiel des candidats d'extraction vers
Mathlib décrits dans \texttt{MATHLIB\_PR\_PLAN.md}.

\paragraph{Approches abandonnées ou révisées.} La première tentative de preuve de
l'indépendance sur la courbe des moments utilisait une matrice de Vandermonde
rectangulaire; elle a été remplacée par l'interpolation polynomiale, plus directe. Le
théorème sur les sous-espaces avait d'abord été prévu avec des séparations explicites
pour $a=0$, $b=0$ et $m=0$, devenues inutiles grâce à l'algèbre extérieure. Enfin, le
dénombrement des événements de séparation avait initialement été formulé par une
division dans les naturels, puis remplacé par une identité multiplicative exacte afin
d'éviter les obligations de divisibilité.

\section{Travaux liés}\label{sec:related}

Les sources mathématiques originales sont Bollobás~\cite{bollobas1965} pour les paires
d'ensembles pondérées et uniformes, Lovász~\cite{lovasz1977} pour la méthode par algèbre
extérieure sur les sous-espaces, et Frankl~\cite{frankl1982} ainsi que
Kalai~\cite{kalai1984} pour les extensions asymétriques. Les notes de
Babai--Frankl~\cite{babaifrankl1992} en donnent une exposition. Mathlib
~\cite{mathlib2020} fournit les puissances extérieures, l'algèbre extérieure graduée, le
déterminant de Vandermonde et l'inégalité LYM. Lors de notre recherche documentée, nous
n'y avons pas trouvé l'inégalité de Bollobás pour les paires d'ensembles, les versions
uniforme ou asymétrique du théorème des deux familles, ni le théorème de Lovász pour les
sous-espaces.

Le code local de Mathlib a été recherché directement. Les index et bibliothèques
externes ont été examinés dans la mesure permise par l'environnement d'audit; une
nouvelle vérification en ligne reste recommandée avant toute affirmation d'originalité
plus forte. Nous n'avons pas localisé de formalisation antérieure réunissant dans un
même développement la preuve pondérée par permutations, le théorème de
Lovász--Frankl pour les sous-espaces et la réduction par courbe des moments vers le
théorème asymétrique de Frankl--Kalai. Nous ne prétendons pas pour autant qu'il s'agit de
la « première formalisation » de chacun de ces résultats pris isolément.

\section{Leçons et limites}\label{sec:lim}

\paragraph{Leçons.} (i) Dans un assistant de preuve, les identités entières exactes sont
souvent préférables aux formulations probabilistes, aussi bien pour la structure de la
preuve que pour les tests. (ii) Une mise en place algébrique uniforme peut absorber les
cas dégénérés au lieu de multiplier les séparations de cas. (iii) La position générale
rectangulaire doit être réellement démontrée, ici par interpolation, et non déduite par
analogie avec la matrice carrée de Vandermonde. (iv) Les transports entre sous-types,
cardinaux et ensembles de positions constituent une part importante du coût de la
preuve combinatoire; ils sont gérés par des isomorphismes d'ordre explicites.

\paragraph{Limites.} Le développement utilise le choix classique et suppose les types
de base et d'indices finis. Il ne formalise pas la classification complète des cas
d'égalité, mais seulement une construction par complémentaires atteignant la borne. La
couche de courbe des moments est formulée sur $\QQ$ plutôt que sur un corps infini
arbitraire. L'ensemble du projet ayant bénéficié d'une assistance importante de systèmes d'IA---notamment pour le cadrage, les preuves, le code, les tests, la rédaction, la traduction, les audits et la préparation du dépôt---,
la vérification par le noyau garantit la correction \emph{par rapport aux énoncés
formels}, mais ne suffit pas à garantir leur fidélité sémantique. Celle-ci est examinée
séparément dans \texttt{INDEPENDENT\_PROOF\_AUDIT.md} et
\texttt{SEMANTIC\_AUDIT.md}; toutes les affirmations restent sous la responsabilité de
l'auteur humain, conformément à \texttt{AI\_USE\_STATEMENT.md}.

\section{Conclusion}\label{sec:concl}

Nous avons formalisé dans Lean~4, selon deux techniques de preuve distinctes, la famille
de théorèmes de Bollobás--Lovász--Frankl sur les deux familles : les versions pondérée
et uniforme de Bollobás par dénombrement de permutations, puis les théorèmes de
Lovász--Frankl pour les sous-espaces et de Frankl--Kalai pour les ensembles asymétriques
par puissances extérieures et position générale de la courbe des moments. Le
développement contient aussi une infrastructure algébrique et combinatoire réutilisable.
Les cinq théorèmes principaux sont vérifiés par le noyau sous les axiomes standards. Au-
delà des résultats particuliers, la formalisation isole une structure commune de
certificat triangulaire derrière deux preuves en apparence éloignées, et organise
plusieurs composants dans une forme proche de celle attendue pour une future intégration
à Mathlib. Ce travail ne revendique aucun nouveau résultat mathématique.

\section*{Disponibilité de l'artefact}
Les sources Lean complètes, les scripts, les documents d'audit, les sources LaTeX et les
versions française et anglaise de ce rapport sont fournis dans l'archive du projet. Le
développement est épinglé sur Lean \texttt{v4.28.0} et Mathlib
\texttt{v4.28.0}; la commande principale de compilation est
\lstinline{lake build}.

\bibliographystyle{plain}
\bibliography{references}

\end{document}
